\scanpage{13}
In the construction of a local Picard scheme $Q$, we have seemed to u\rlap{n}se in a technically essential way the \tikz[baseline=(s.base)]{\node[inner sep=0pt](s){\deleted{fact}};\node[anchor=south west,inner sep=1pt](a) at ([xshift=-.5em,yshift=7pt]s.north west){\inserted{assumption}};} that $U$ is regular (plus resolution). I would conjecture that a reasonable \tikz[baseline=(s.base)]{\node[inner sep=0pt](s){loc al};\draw[line width=.3pt,gray] ([xshift=-1pt,yshift=-2pt]s.center) to[out=-10,in=200] ([xshift=6pt,yshift=1pt]s.center);} Picard scheme $Q$, \tikz[baseline=(s.base)]{\node[inner sep=0pt](s){can};\node[anchor=south west,inner sep=1pt](a) at ([xshift=-1em,yshift=7pt]s.north west){\inserted{locally of finite type over $k$,}};\draw ([xshift=-3pt,yshift=-2pt]s.south west) to[out=90,in=180] ([xshift=0pt,yshift=-2pt]a.south west)--([xshift=0pt,yshift=-2pt]a.south east) to[out=0,in=-90] ([xshift=3pt,yshift=3pt]a.south east);} be constructed also without this assumption (provided that a field of representatives is given, \rlap{l}or the residue field is perfect of char. $p\handwritten{>}0$, of course). Instead of a resolution of singularities, one \rlap{\deleted{p}}\handwritten{m}ay then wish to find a proper surjective morphism
\[
f:X\handwritten{\longrightarrow}S
\]
inducing an isomorphism $f^{-1}(U)\handwritten{\xrightarrow{\sim}}U$, and such that the corresponding map
\[
\operatorname{Pic}(X)\handwritten{\longrightarrow}\operatorname{Pic}(U)
\]
is \uline{surjective}. We then could repeat the construction of $P$, as $\lim\underline{\operatorname{Pic}}_{X_n/k}$, as above, and of $Q$ as $\operatorname{Coker}(\underline{D}\handwritten{\longrightarrow}P)$, where $\underline{D}$ is a group scheme corresponding to divisors concentrated on $X_o$ (possibly no longer discrete ?). Another way of approach would be to take a resolution of singularities $f:X\handwritten{\longrightarrow}S$, but where now we cannot assume any longer that $f^{-1}(U)\handwritten{\longrightarrow}U$ is an isomorphism, and to describ\rlap{x}e\deleted{\unsure{s}} a scheme $P$ in terms of invertible sheaves on $X$, together with \uline{descent data} of $\underline{L}\handwritten{|}f^{-1}(U)$ to $U$. A nice test for the yoga of local Picard schemes without regularity assumptions on $U$ would be to see if it is true in general that \rlap{f}\kern.2em or $A$ Cohen-Macaulay of dim $\handwritten{\geq3}$ (more generally, for $H^2_{\underline{m}}(A)=0$), the group $\operatorname{Pic}(U)$ is indeed countable, and for separably closed residue field \tikz[baseline=(s.base)]{\node[inner sep=0pt](s){\deleted{\unsure{invariant}}};\node[anchor=south west,inner sep=1pt](a) at ([xshift=0em,yshift=7pt]s.north west){\inserted{is invariant under}};} \uline{separable} field extension $k\handwritten{\longrightarrow}k'$, $A$ being replaced by $A\handwritten{\widehat{\mathord{\tikz[baseline=(s.base)]{\node[inner sep=0pt](s){$\mathrm{x}$};\draw[line width=.5pt] ([xshift=.5pt,yshift=1.3pt]s.center) ellipse[x radius=3pt,y radius=4.5pt];}}}}_k k'$ (which would express the expected \unsure{discrete|discreto} structure of the local Picard scheme).
% archive sheet number: 12
